\documentclass[10pt]{article}
\usepackage{template/template}
\usepackage{tikz}
\NewsSetKeys{date={2026年11月8日},vol={026},title={每日有趣新闻},author={编辑部},note={待验收短刊｜新闻核查截至北京时间11月8日上午。本期尚未完成完整深稿、正式图像归档与PDF验收，不应视为正式发行。}}
\begin{document}
\begin{documents}
\newsitem{一场首都里的对立集会：印尼经济与社会政策承受双重压力}{国际政治与社会｜11月7日｜路透社}{
11月7日，印度尼西亚首都雅加达出现两组政治诉求相反的集会：数百名工人表达对总统普拉博沃的支持；另有数百名妇女及民间社会参与者批评政府的妇女权利政策、免费供餐项目和卫生用品税收安排。两场活动不应简单合并为一场示威。前者着眼于政府支持和社会稳定，后者强调公共政策对不同群体的影响。路透社报道还指出，野火应对与投资者信心下降构成政府面临的背景压力。\cite{reutersindonesia}\par
与此同时，警方将四名高级官员列为涉嫌试图监听总统通信的嫌疑人，其中一人与前总统佐科·维多多有关联。嫌疑人与被定罪者不同；政治关系也不能代替司法证据。前总统表示尊重法律程序，相关律师称当事人将配合调查。判断事件是否意味着政治联盟破裂，仍需后续司法材料和各方独立表态。\cite{reutersindonesia}\par
编辑观察：同一日支持与反对集会并存，显示社会诉求不能用单一民意标签概括。集会人数也不能直接推算全国支持率。}
\newsitem{霍尔木兹海峡谈判：调解进展不是通航协议}{战争与外交｜11月7日｜路透社}{
卡塔尔首相11月7日表示，美国与伊朗之间的调解仍在继续，并认为达成协议存在可能性。这属于调解方的公开评估，并不意味着美伊双方已经签署停火或恢复海峡航行的协议。报道提到，伊朗提出的方案包括解除港口封锁、解冻资产以及海峡通航安排，但美方此前未接受该方案。有关格什姆岛疑似遭导弹袭击的消息来自伊朗半官方媒体的初步报道，尚无足够官方证据确认损害情况，更不能由此断言袭击责任方。\cite{reutersqatar}\par
对于依赖海上能源运输的国家而言，真正值得跟踪的是航道保险、船舶通行、军事行动与谈判执行之间的差别。外交官对协议的乐观表态与商船恢复通航，属于两类不同的可核验事实。}
\newsitem{科技股下跌之后，央行担忧为何从通胀转向增长}{经济与AI｜11月6—7日｜路透社}{
本周科技股持续回调，市场对数据中心投资的回报、融资和监管前景出现新的疑虑。路透社分析指出，股票估值下降可能通过财富效应和融资条件影响消费与投资；与此同时，中东战争引起的能源成本上涨仍对物价形成压力。央行面对的并非简单的加息或降息选择，而是增长减速与通胀偏高可能同时出现的局面。\cite{reutersmarkets}\par
另一则11月6日就业报道显示，美国10月新增就业16.5万人、失业率降至4.1\%，市场对12月加息的押注上升。但期货定价是市场预期，不是联储决议；不同时间截面的交易价格也可能迅速改变。\cite{reutersfed}\par
在人工智能投资问题上，需要区分实际已经建设的数据中心、企业宣布的未来资本开支与股票市场赋予的增长估值。短期市值损失不等于相同金额的实体资本消失。}
\newsitem{两艘火星探测器为何要分别借地球加速？}{航天｜11月7日｜NASA任务说明}{
NASA的ESCAPADE任务由两艘探测器组成，计划分别在11月7日和9日近距离飞掠地球，同时进行约九分钟的发动机点火，获得飞往火星所需的额外速度。与单纯的重力弹弓不同，这种动力飞掠将轨道几何与推进系统的速度增量结合。NASA在11月2日的预告称，两艘探测器预计于2027年9月抵达火星，随后协同研究太阳风与火星磁环境的相互作用及大气逃逸。\cite{nasaescapade}\par
本刊截稿时仅确认NASA事前任务说明，尚未取得11月7日飞掠完成后的独立任务状态通报，因此不将计划时刻误写为已经成功完成的飞掠。}
\newsitem{月球南极的山峰：为何光照和阴影都重要？}{航天影像｜11月7日｜NASA APOD}{
NASA11月7日发布Alejandro Sanz于10月31日拍摄的月缘影像，展示接近月球南极的De Gerlache–Kocher山体。其高出周围约四公里，附近既有可能长期保存水冰的永久阴影坑底，也有较长时间可获得阳光的高地。影像的斜视透视受到当时月球天平动角度影响，不能把照片中的视距当作精确地形比例尺。\cite{nasamoon}\par
图片原始来源已核实，但摄影者署名并不自动授予本刊再发布许可；本次没有将该照片复制进仓库。}
\newsitem{数学检索与验证状态}{数学｜11月6—8日}{
已核查近期arXiv数学与Lean相关检索结果，但当前可访问的结果未能可靠确认11月6—7日新提交论文的作者、版本和核心定理。因此本期不发布所谓“最新证明”。对较早的AI生成数学稿件综述，作者明确说明其新颖性判断暂定，且没有对相关Lean开发进行内核编译验证；它不构成那些原始证明正确性的独立认证。\cite{arxivsurvey}}
\newsitem{出版说明}{编辑部}{
本期为经过部分来源核查的待验收短刊。尚缺少足够长度的深稿、可合法归档的实际图像、GitHub Actions构建日志及正式PDF逐页验收。不得以短刊源文件代替已发行的完整报纸。}
\bibliographystyle{plain}\bibliography{refs}
\end{documents}
\end{document}